

%%%PACKAGES
\usepackage[utf8]{inputenc}
\usepackage[T1]{fontenc}
\usepackage[french]{babel}
\usepackage{multicol}
\usepackage[left=2.1cm, right=2.1cm, bottom=2.7cm, top=2.7cm]{geometry}
\usepackage{titling}
\usepackage{xcolor}
\usepackage{fancyhdr}
\pagestyle{fancy}
\usepackage{fancybox}
\usepackage{lastpage}
\usepackage{stmaryrd}
\usepackage[normalem]{ulem}
\usepackage{graphicx}
\usepackage{soul}
\usepackage{float}
\usepackage{amsmath,amsfonts,amsthm,amssymb} 
\usepackage{mdframed}
\usepackage{enumitem}
\usepackage{changepage}
\usepackage{titlesec}
\usepackage[most]{tcolorbox}
\usepackage{tabto,lipsum}
\usepackage{marginnote}[] %%%pour les symboles dans la marge
\usepackage{tasks} % pour aligner des enumerate
\usepackage{framed}

%%%PACKAGE SYMBOLES
\usepackage{pifont}
\usepackage{fourier}
\usepackage{ticollege} %pour le symbole calculatrice

%%%PACKAGES TIKZ
\usepackage{tikz}
\usepackage{pstricks}
\usepackage{pstricks-add}
\usetikzlibrary{arrows}
\usepackage{mathrsfs}
\usepackage{pgfplots}
\pgfplotsset{compat=1.15}




%%%RACCOURCIS
\newcommand{\N}{\mathbb{N}}
\newcommand{\Z}{\mathbb{Z}}
\newcommand{\Q}{\mathbb{Q}}
\newcommand{\R}{\mathbb{R}}
\newcommand{\C}{\mathbb{C}}
\newcommand{\oa}{\overrightarrow}
\newcommand{\ol}{\overline}


%%% en haut et en bas
\fancyhf{}
\fancyhead[L]{\type}
\fancyhead[R]{\niveau}
\fancyfoot[L]{Ch.\numchap : \titrechap}
\fancyfoot[R]{\thepage \hspace{1pt} / \pageref{LastPage}}


%%%ENCADRÉS
%énoncés importants
\newcommand{\cbox}[1]{\begin{tcolorbox}[lower separated=false, colback=white]
#1
\end{tcolorbox}}
%%%LEFTBAR
\renewenvironment{leftbar}{%
  \def\FrameCommand{\hspace{1cm}\vrule width 0.4pt \hspace{10pt}}%
  \MakeFramed{\advance\hsize -\width\FrameRestore}}%
 {\endMakeFramed}



\setlength\parindent{0pt}


%%% STYLE DES SECTIONS
\renewcommand\thesection{\Roman{section}.}
\renewcommand\thesubsection{\arabic{subsection}.}
\renewcommand\thesubsubsection{\alph{subsubsection})}

\titleformat{\section}
  {\scshape \Large \bfseries}{\thesection}{1em}{}
\titleformat{\subsection}
  {\normalfont  \large \bfseries}{\thesubsection}{1em}{}
\titleformat{\subsubsection}
  {\normalfont\large }{\thesubsubsection}{1em}{}

%%%STYLE DES ITEMIZE
\setitemize[1]{label=\textbullet}


%%%THEOREM STYLES

\newtheoremstyle{exercices}
{}                % Space above
{}                % Space below
{}        % Theorem body font % (default is "\upshape")
{}                % Indent amount
{\bfseries}       % Theorem head font % (default is \mdseries)
{}               % Punctuation after theorem head % default: no punctuation
{10pt}               % Space after theorem head
{}                % Theorem head spec
\theoremstyle{exercices}
\newtheorem{exo}{Exercice}

\newtheoremstyle{enonce}
{}                % Space above
{}                % Space below
{}        % Theorem body font % (default is "\upshape")
{}                % Indent amount
{\bfseries}       % Theorem head font % (default is \mdseries)
{ :}               % Punctuation after theorem head % default: no punctuation
{3pt}               % Space after theorem head
{}                % Theorem head spec
\theoremstyle{enonce}
\newtheorem*{defi}{Définition}
\newtheorem*{defis}{Définitions}
\newtheorem*{prop}{Propriété}
\newtheorem*{props}{Propriétés}
\newtheorem*{thm}{Théorème}
\newtheorem*{defiprop}{Définition-propriété}

\theoremstyle{remark}
\newtheorem*{demo}{Démonstration}
\newtheorem*{exple}{Exemple}
\newtheorem*{exples}{Exemples}
\newtheorem*{rem}{Remarque}
\newtheorem*{rems}{Remarques}
\newtheorem*{note}{Notation}
\newtheorem*{met}{Méthode}
\newtheorem*{app}{Application}
\newtheorem*{ill}{Illustration géométrique}

%%%TITRE
\newcommand{\montitre}[1]{
\begin{center}
%\vspace{-0.5cm}
\doublebox{
 \Large \scshape 
 #1}
\end{center}
}


%%%CITATION
\renewcommand{\citation}[4]{
\begin{minipage}{0.25\textwidth}
\hspace{1pt}
\end{minipage} \hfill
\begin{minipage}{0.7\textwidth}
\center
\small
\og #1 \fg{}
\flushright {\footnotesize \textsc{#2}, \textit{#3} (#4)}
\end{minipage}\\[\baselineskip]
}


%%%SYMBOLES MARGE
\newcommand{\acomp}{\reversemarginpar
\marginnote{\leafNW}}

\newcommand{\attention}{\reversemarginpar
\marginnote{\warning}}

\newcommand{\PCQCI}{\reversemarginpar
\marginnote{\ding{106}}}

\newcommand{\todo}{\reversemarginpar
\marginnote{\ding{217}}}

\newcommand{\calc}{\reversemarginpar
\marginnote{\TiCCalc*[calcscale=0.4, calcrotate=0]}}

